\documentclass[11pt,a4paper]{article}
\usepackage[T1]{fontenc}
\usepackage[utf8]{inputenc}
\usepackage{lmodern}
\usepackage[margin=25mm]{geometry}
\usepackage{amsmath,amssymb,booktabs,longtable,array,tabularx}
\usepackage{hyperref}
\usepackage{fancyvrb}
\usepackage{microtype}
\usepackage{graphicx}
\setlength{\emergencystretch}{3em}
\hypersetup{colorlinks=true,linkcolor=blue,citecolor=blue,urlcolor=blue}
\renewcommand{\refname}{Bibliography}
\title{Exception-rules\\Scientific background and experimental evaluation\\\large Supplementary materials}
\author{Dawid Macha \and Łukasz Wróbel \and Marek Sikora}
\date{October 2026}
\begin{document}
\maketitle
\begin{abstract}
This supplement presents the theoretical foundations, induction algorithms, and experimental evaluation of the exception-rules package for classification, regression, and survival analysis. The evaluation covers exception discovery on 49 classification, 48 regression, and 25 survival datasets, illustrated by explicit rule structures and summaries. Computational results compare fitting times with conventional rule learners and decision trees and describe their association with dataset dimensions.
\end{abstract}
\setcounter{tocdepth}{2}
\tableofcontents
\clearpage



% Editable supplement chapter.
\section{Theoretical basics}
\label{supplement-theory}




\subsection{Definition of exception rules}
\label{suzuki-s-formulation-and-the-reference-rule}

The classical formulation attributed to Suzuki and Shimura
\cite{Suzuki1996} uses two rules:

\[
\mathrm{CR}: A\to X, \qquad
\mathrm{ER}: A\land B\to Y, \quad Y\ne X.
\]

The commonsense rule (CR) represents a dominant pattern with high
support and confidence. The exception rule (ER) introduces an
additional condition \(B\), identifying a smaller subgroup in which
the outcome differs from the general conclusion. Consequently,
exception rules typically have lower support but high confidence
in the alternative class.

For a classification rule, define

\[
\operatorname{supp}(A\to X)
  =\frac{|C_A|}{n},\qquad
\operatorname{conf}(A\to X)
  =\frac{|\{i\in C_A:y_i=X\}|}{|C_A|}.
\]

Here support measures the proportion of observations satisfying
the premise. Class coverage, \(p/P\), is the proportion of all
observations of class \(X\) covered by the rule.
Confidence is also called precision.

However, the CR--ER pair alone does not guarantee that the
exception represents a non-trivial dependency. In particular,
if \(B\) already strongly predicts \(Y\), the exceptional
outcome may be explained by \(B\) independently of \(A\).
In such a case, adding \(A\) does not necessarily reveal
a new relationship.

To address this limitation, Hussain et al. \cite{Hussain2000}
introduced a \textbf{reference rule (RR)}:

\[
\mathrm{RR}: B\to Y.
\]

The reference rule evaluates the relationship between the
additional condition \(B\) and the exceptional outcome \(Y\)
without requiring \(A\). It provides a baseline against which
the exception rule can be assessed. Ideally, this relationship
has low support and/or low confidence, indicating that \(B\)
alone does not adequately explain the exceptional outcome.
Conversely, if \(B\to Y\) already has high confidence, the
exception may simply reflect a strong independent association
between \(B\) and \(Y\), rather than a distinctive effect
of combining \(A\) and \(B\).

A stronger, interaction-oriented interpretation arises when
\(B\) independently supports the same outcome as the
commonsense rule. This leads to the following formulation,
adopted by the task-specific methods considered here:

\[
\mathrm{CR}: A\to X,\qquad
\mathrm{RR}: B\to X,\qquad
\mathrm{ER}: A\land B\to Y,\quad Y\ne X.
\]

In this setting, both \(A\) and \(B\) separately favor the
general outcome \(X\), whereas their conjunction favors
the alternative outcome \(Y\). The reference rule therefore
serves a more specific purpose: it establishes that the
additional condition \(B\), considered independently,
supports the general pattern rather than the exception.

This makes the exception particularly informative. Neither
parent premise alone suggests the exceptional outcome,
yet their intersection exhibits a different dominant class.
The resulting reversal cannot be explained by the marginal
association of \(B\) with \(Y\) alone and instead indicates
a context-dependent relationship between the premises.

The distinction between these two formulations is important.
The original reference rule \(B\to Y\) assesses whether the
exceptional outcome is already associated with \(B\).
In contrast, the stronger formulation \(B\to X\) explicitly
requires both parent rules to favor the general outcome,
making the reversal within their intersection the central
pattern of interest.





\subsection{Classification exception rules}
\label{classification-exception-rules}
In classification, \(X\) and \(Y\) are class labels. The adopted
structure requires strong CR and RR rules predicting \(X\) and a strong
intersection predicting a different class \(Y\). The precision threshold
\(\theta\) controls these requirements. In a multiclass problem,
\(Y\ne X\) is a specific alternative label, rather than a single merged
class containing all other labels. The implementation chooses the most frequent
alternative class in the intersection.



\subsection{Regression exception rules}
\label{regression-exception-rules}

The concept of exception rules can be naturally extended from classification to regression by replacing discrete class labels with continuous target values. While in classification the objective is to identify examples with different class assignments, in regression the goal is to detect subgroups whose expected values differ from the general pattern.

Analogously to the classification setting, regression exception rules consist of three components: a \emph{commonsense rule (CR)}, a \emph{reference rule (RR)}, and an \emph{exception rule (ER)}:

\begin{center}
\[
\begin{aligned}
\text{CR: } & A \rightarrow \mu_{A} \; [l_{A}, u_{A}], \\
\text{ER: } & A \land B \rightarrow \mu_{AB} \; [l_{AB}, u_{AB}], \\
\text{RR: } & B \rightarrow \mu_{B} \; [l_{B}, u_{B}].
\end{aligned}
\]
\end{center}

Each rule is associated with a numerical estimate and a descriptive variability band: the covered-target mean $\pm$ one standard deviation. Its role here is an operational search criterion, not an uncertainty bound on the expected outcome.

For a rule covering a set of examples $S$, the predicted value and its variability interval are defined as
\[
\mu_S = \frac{1}{|S|}\sum_{i \in S} y_i, \qquad
l_S = \mu_S - \sigma_S, \qquad
u_S = \mu_S + \sigma_S,
\]
where $\sigma_S$ denotes the standard deviation of the target variable for the examples in $S$. \cite{Janssen2011} \cite{Sikora2012}

The commonsense rule captures the dominant trend in the data, while the reference rule verifies whether the additional conditions alone influence the target variable. The exception rule identifies a subgroup whose expected value deviates from the general pattern.

A regression exception rule is defined by the following conditions:

\begin{itemize}
\item Consistency of the commonsense and reference rules:
\[
\mu_{A} \in [l_{B}, u_{B}] \quad \text{and} \quad
\mu_{B} \in [l_{A}, u_{A}].
\]

\item Deviation of the exception rule from the commonsense rule:
\[
\mu_{AB} \notin [l_{A}, u_{A}].
\]

\item Deviation of the exception rule from the reference rule:
\[
\mu_{AB} \notin [l_{B}, u_{B}].
\]
\end{itemize}

These conditions define a candidate whose parent means are compatible with the specified bands and whose intersection mean lies outside both bands. They do not establish equal parent expectations, statistical significance, causality, or practical importance. The resulting pattern is a hypothesis about context-dependent association.

Candidates are additionally screened using coverage and Mann--Whitney U comparisons \cite{Mann1947}. Screening on the discovery data is not independent validation and does not control the false-discovery rate of the full search. The precise groups and acceptance thresholds are given in Section~\ref{algorithm-regression}.

\subsection{Survival exception rules}
\label{survival-exception-rules}
For a covered examples with event times $t_j$, event counts $d_j$, and
numbers at risk $n_j$ immediately before each event time, its Kaplan--Meier
conclusion is \cite{MachaSurvival2026}
\[
\widehat S(t)=\prod_{t_j\leq t}(1-d_j/n_j).
\]
Right-censored observations contribute to the risk set until censoring.
The median is the first time at which $\widehat S(t)\leq0.5$; if that level
is not reached, the median is not estimable within follow-up. Restricted
mean survival time is $\int_0^\tau\widehat S(t)\,dt$ for a stated horizon $\tau$.

The concept of exception rules can be naturally extended from classification to survival analysis by replacing discrete class labels with survival distributions. While in classification the goal is to identify subgroups with different class assignments, in survival analysis the objective is to identify subgroups whose time-to-event distributions differ significantly from the general pattern.

Analogously to the classification setting, survival exception rules consist of three components: a \emph{commonsense rule (CR)}, a \emph{reference rule (RR)}, and an \emph{exception rule (ER)}:

\begin{center}
\(\text{CR: } A \;\rightarrow\; S_X(t)\)

\(\text{ER: } A \land B \;\rightarrow\; S_Y(t)\)

\(\text{RR: } B \;\rightarrow\; S_{X'}(t)\)
\end{center}

Each rule is associated with an empirical survival function estimated from
the observations satisfying its premise. CR defines the parent subgroup
with estimate $S_X(t)$, RR defines the second parent with estimate $S_{X'}(t)$,
and ER defines their intersection with estimate $S_Y(t)$. The RR comparison
screens for candidates whose parent curves are not separated by the chosen
test; it does not establish that their survival distributions are equivalent.

%The key requirement is defined in terms of statistical relationships between survival distributions:
%\begin{itemize}
%\item no statistically significant difference is detected between the survival distributions of CR and RR,
%\item statistically significant differences are observed between the survival distributions of CR and ER,
%\item statistically significant differences are observed between the survival distributions of RR and ER.
%\end{itemize}

The operational screening requirement concerns comparisons among three groups -- CR, RR, and ER:

\begin{itemize}
\item CR vs RR: no statistically significant difference,
\item CR vs ER: statistically significant difference detected,
\item RR vs ER: statistically significant difference detected.
\end{itemize}

These relationships are typically assessed using statistical hypothesis testing procedures, such as the log-rank test. Since the subgroups defined by the rules may overlap (e.g., the ER subgroup is a subset of both CR and RR), these comparisons should be interpreted as approximate measures of differences in survival behavior rather than strictly independent group comparisons.




% Editable supplement chapter.
\section{Algorithms}
\label{supplement-algorithms}

\subsection{Shared induction strategy}
\label{algorithms}

\label{shared-induction-strategy}
All three learners use separate-and-conquer (sequential covering). They grow a
conjunctive rule greedily, retain it if it covers new observations, remove its
coverage from the working uncovered set, and repeat. The full training set
remains available for evaluating rule quality. Removing an observation from the
uncovered set does not remove it from subsequent quality calculations.

Candidate conditions test nominal equality or numerical thresholds between
adjacent distinct observed values. Candidate generation is local to the
currently covered data. Classification optionally restricts numerical cuts to
class boundaries. Minimum coverage prevents a new rule from simply repeating
previous coverage; {\ttfamily\small max\_growing} limits CR and RR premise lengths separately.
Consequently an ER conjunction may contain more conditions than either parent.

The pseudocodes below describe the methods. They expose the search and acceptance
criteria, rather than serving as line-by-line replacements for the Python
implementation. 
Here {\ttfamily\small Cov(A)} denotes premise coverage in the full training set, {\ttfamily\small U} is the
working uncovered set, and an empty premise initially covers all observations.

\noindent\begin{minipage}{\linewidth}
\textbf{Algorithm 1. Sequential covering (common structure).}


\begin{Verbatim}[fontsize=\scriptsize,frame=single,samepage=true]
INDUCE(D, task, mincov, max_growing):
    R <- empty list
    targets <- class labels for classification; one task for other problems
    for target in targets:
        U <- target-class observations for classification; all D otherwise
        while U is not empty:
            (CR, RR, ER) <- GROW_WITH_EXCEPTION_SEARCH(D, U, target)
            if CR has no conditions: break
            apply task-specific finalisation (see pruning policy below)
            newly_covered <- eligible coverage of CR intersect U
            if newly_covered is empty: break
            append (CR, RR, ER) to R
            U <- U minus newly_covered
    update rule conclusions, coverage statistics, and model metadata
    return R
\end{Verbatim}
\end{minipage}
Classification removes covered positives of the current class from {\ttfamily\small U}.
Regression and survival remove covered observations. A rule without an accepted
exception is still a valid output. Candidate coverage checks also allow final
small remainders in the implementation; {\ttfamily\small mincov} is not a global lower bound
on every reported subgroup size.



\subsection{Classification}
\label{classification-ceris}
The operational classification constraints are
\[
\operatorname{conf}(A\to X)>\theta,\qquad
\operatorname{conf}(B\to X)>\theta,\qquad
\operatorname{conf}(A\land B\to Y)>\theta,\quad Y\ne X.
\]
guarantee in the reported experiments.
The CR predicts a fixed target class \(X\) and grows according to the
configured quality measure. If it is sufficiently precise but still covers
other classes, the learner searches for an RR predicting the same class.
The intersection may predict another class with high precision.

\noindent\begin{minipage}{\linewidth}
\textbf{Algorithm 2. Classification commonsense-rule growth.}


\begin{Verbatim}[fontsize=\scriptsize,frame=single,samepage=true]
GROW_CLASSIFICATION(D, U, X):
    A <- empty premise; RR <- none; ER <- none
    while the growth limit has not been reached:
        c <- highest-quality admissible condition extending A
        if c does not exist: break
        A <- A AND c
        if A covers no observations outside class X: break
        if precision(A -> X) > theta and exception search is enabled:
            (RR, ER) <- SEARCH_CLASSIFICATION_REFERENCE(A, D, X)
            if ER exists: break
    finalise the CR premise using the task's quality history
    return (A -> X, RR, ER)
\end{Verbatim}
\end{minipage}
RR search must simultaneously support \(X\) outside the CR and retain some
of the CR's negatives so that the conjunction can express a class reversal. The current implementation
uses the normalised score


\[
s(B)=q(B\to X)
   +\frac{|C_B\cap U_X|}{|U_X|}
   +\frac{|C_B\cap N_A|}{|N_A|},
\]
where \(N_A=\{i\in C_A:y_i\ne X\}\) and \(U_X\) is the reference
search's set of uncovered positives outside the CR. A ratio with an empty
denominator contributes zero. 

\noindent\begin{minipage}{\linewidth}
\textbf{Algorithm 3. Classification reference and exception search.}


\begin{Verbatim}[fontsize=\scriptsize,frame=single,samepage=true]
SEARCH_CLASSIFICATION_REFERENCE(A, D, X):
    B <- empty premise
    U_X <- class-X observations outside A
    N_A <- non-X observations covered by A
    repeat until growth limit or no eligible refinement:
        evaluate candidate extensions B' using reference score s(B')
        require new positive coverage and precision(B' -> X) > theta
        test the intersection A AND B' for an alternative class Y
        prefer valid exception candidates using the selector's score rules
        extend B with the selected condition
        if the selected intersection predicts Y != X with precision > theta:
            return (B -> X, A AND B -> Y)
    return (none, none)
\end{Verbatim}
\end{minipage}
The alternative label is the modal class among the intersection's non-X
observations, and its precision is calculated on the entire intersection.
The selector can favour an admissible exception candidate before it improves
the previous score; it is not simply a global maximisation of \(s\).
An empty intersection of negatives cannot produce an attached exception.


\subsection{Regression}
\label{algorithm-regression}
To construct rule positive/negative counts,
the learner calculates the covered target mean and standard deviation. Targets
inside this local mean-plus/minus-standard-deviation band are treated as
positive, and the same band defines dataset-wide positive totals. This is a
quality calculation for numerical targets, rather than conversion to fixed
class labels.

For the statistical screen, let $E=C_A\cap C_B$, $G_A=C_A\setminus E$,
and $G_B=C_B\setminus E$. The three groups are disjoint. The current code
accepts a candidate only when the two-sided Mann--Whitney comparisons satisfy
\[
p(G_A,G_B)>0.05,\qquad p(E,G_A)\leq0.05,\qquad p(E,G_B)\leq0.05.
\]


\noindent\begin{minipage}{\linewidth}
\textbf{Algorithm 4. Growing a regression commonsense rule.}


\begin{Verbatim}[fontsize=\scriptsize,frame=single,samepage=true]
GROW_REGRESSION(D, U):
    A <- empty premise
    while the growth limit has not been reached:
        c <- highest-C2 admissible condition extending A
        if c does not exist: break
        A <- A AND c
        update covered-target mean and variability band
        if exception search is enabled:
            (RR, ER) <- SEARCH_REGRESSION_REFERENCE(A, D)
            if ER exists: return (A, RR, ER)
    return (A, none, none)
\end{Verbatim}
\end{minipage}
\noindent\begin{minipage}{\linewidth}
\textbf{Algorithm 5. Regression reference search.}


\begin{Verbatim}[fontsize=\scriptsize,frame=single,samepage=true]
SEARCH_REGRESSION_REFERENCE(A, D):
    B <- empty premise; best <- none; best_overlap <- 0
    repeat until growth limit:
        selected <- none
        for each unused elementary condition c extending B:
            B' <- B AND c; E' <- A AND B'
            require sufficient new reference coverage outside A
            require compatible parent means and exceptional intersection mean
            require the three disjoint-group Mann-Whitney criteria
            if |Cov(E')| > best_overlap:
                selected <- c; best <- (B', E')
                best_overlap <- |Cov(E')|
        if selected is none: break
        B <- B AND selected
    return best, or (none, none)
\end{Verbatim}
\end{minipage}

\subsection{Survival analysis}
\label{algorithm-survival-analysis}
Survival CR conditions maximise log-rank separation between covered and
uncovered observations. RR candidates are feasible if the CR and RR comparison
does not reject at \(0.05\), while both comparisons involving their
intersection reject at that level. Among feasible candidates, the reference
selector favours a larger CR--RR p-value. This ranking heuristic is not an
equivalence test.

\noindent\begin{minipage}{\linewidth}
\textbf{Algorithm 6. Survival CR and reference search.}


\begin{Verbatim}[fontsize=\scriptsize,frame=single,samepage=true]
GROW_SURVIVAL(D, U):
    A <- empty premise
    repeat until growth limit:
        c <- admissible condition with highest covered-vs-uncovered log-rank
        if c does not exist: break
        A <- A AND c
        if exception search is disabled: continue
        B <- empty premise; best_score <- 0
        repeat until reference growth limit:
            for candidate B' = B AND c_ref:
                E' <- A AND B'
                fit Kaplan-Meier curves on Cov(A), Cov(B'), Cov(E')
                calculate p_AB, p_AE, p_BE by log-rank comparisons
                require nonempty E' and sufficient new coverage outside A
                require p_AB > 0.05 and p_AE <= 0.05 and p_BE <= 0.05
                retain candidate if p_AB improves best_score
            if no candidate retained: break
            B <- retained candidate; update best_score
        finalise B using the implementation's quality history
        recheck intersection separation before attaching RR and ER
        if accepted: break
    finalise A and fit its Kaplan-Meier conclusion
    return (A, attached RR or none, attached ER or none)
\end{Verbatim}
\end{minipage}
Prediction uses the KM estimator of an applicable exception in place of its
parent rule's estimator, then averages the estimators contributed by covering
rules. The source of each prediction remains
inspectable through the rule structure.





% Editable supplement chapter.
\section{Experiments}
\label{supplement-experiments}

\subsection{Evidence and evaluation questions}
\label{experiments}

\label{evidence-and-evaluation-questions}
The experiments document exception-rule discovery: how many rules were
generated and which premises define the discovered exceptions. Discovery was
performed on complete datasets. Computational cost and timing comparisons are collected in
Section~\ref{supplement-computational-efficiency}.

The evaluation covers classification, regression, and survival analysis.
The discovery results are consolidated from the earlier studies
\cite{MachaRegression2026,MachaSurvival2026};
For each task, dataset-level counts are accompanied by examples of the
discovered rule structures.


\begin{table}[htbp]
\centering
\small
\caption{Overview of the discovery experiments.}
\begin{tabularx}{\linewidth}{@{}lrX@{}}
\toprule
Task & Datasets & Discovery settings \\
\midrule
Classification & 49 & C2; mincov 5 or 0.1; theta 0.6--0.9; growth limit 5 \\
Regression & 48 & C2; mincov 5; growth limit 5 \\
Survival & 25 & Log-rank; mincov 5; growth limit 5 \\
\bottomrule
\end{tabularx}
\end{table}

The experiments use benchmark datasets drawn from public repositories,
with task-specific induction settings.
The minimum new coverage is five for regression and survival. Classification
uses either five observations or 10\% of the target class, combined with
precision thresholds of 0.6, 0.7, 0.8, and 0.9. This grid characterises
parameter sensitivity. The five-observation setting follows the RuleKit
configuration adopted in the regression study \cite{MachaRegression2026}.

C2 is used for classification and regression rule quality
\cite{Wrobel2016}; survival induction uses log-rank separation. The growth
limit is five conditions per parent rule to constrain premise length.
These settings define the experimental protocol rather than a predictive
hyperparameter optimisation. 


\subsection{Classification}
\subsubsection{Discovery across datasets}
\label{classification}

\label{discovery-across-datasets}
Classification exception discovery uses complete datasets.
Fractional {\ttfamily\small mincov=0.1} means a fraction of the
current target class, whereas {\ttfamily\small mincov=5} means an absolute count of new
positive observations. The precision threshold applies to the three components
of the exception structure.


\begingroup
\small
\setlength{\tabcolsep}{4pt}
\begin{longtable}{>{\raggedright\arraybackslash}p{\dimexpr 0.20000\linewidth-2.0000\tabcolsep\relax}>{\raggedright\arraybackslash}p{\dimexpr 0.20000\linewidth-2.0000\tabcolsep\relax}>{\raggedright\arraybackslash}p{\dimexpr 0.20000\linewidth-2.0000\tabcolsep\relax}>{\raggedright\arraybackslash}p{\dimexpr 0.20000\linewidth-2.0000\tabcolsep\relax}>{\raggedright\arraybackslash}p{\dimexpr 0.20000\linewidth-2.0000\tabcolsep\relax}}
\caption{Datasets with at least one exception, out of 49.}\\
\toprule
mincov & theta = 0.6 & theta = 0.7 & theta = 0.8 & theta = 0.9 \\
\midrule
\endfirsthead
\toprule
mincov & theta = 0.6 & theta = 0.7 & theta = 0.8 & theta = 0.9 \\
\midrule
\endhead
5 & 13 & 8 & 4 & 2 \\
0.1 & 6 & 4 & 3 & 2 \\
\bottomrule
\end{longtable}
\endgroup

\begingroup
\small
\setlength{\tabcolsep}{4pt}
\begin{longtable}{>{\raggedright\arraybackslash}p{\dimexpr 0.20000\linewidth-2.0000\tabcolsep\relax}>{\raggedright\arraybackslash}p{\dimexpr 0.20000\linewidth-2.0000\tabcolsep\relax}>{\raggedright\arraybackslash}p{\dimexpr 0.20000\linewidth-2.0000\tabcolsep\relax}>{\raggedright\arraybackslash}p{\dimexpr 0.20000\linewidth-2.0000\tabcolsep\relax}>{\raggedright\arraybackslash}p{\dimexpr 0.20000\linewidth-2.0000\tabcolsep\relax}}
\caption{Total exception rules across the 49 datasets.}\\
\toprule
mincov & theta = 0.6 & theta = 0.7 & theta = 0.8 & theta = 0.9 \\
\midrule
\endfirsthead
\toprule
mincov & theta = 0.6 & theta = 0.7 & theta = 0.8 & theta = 0.9 \\
\midrule
\endhead
5 & 33 & 15 & 6 & 5 \\
0.1 & 13 & 7 & 5 & 3 \\
\bottomrule
\end{longtable}
\endgroup
More restrictive precision requirements yield fewer exceptions in these
experiments. This demonstrates parameter sensitivity; it does not establish
that every retained exception is useful or every rejected one is noise.


\subsubsection{Exception counts by dataset}
\label{classification-counts-by-dataset}
The following table summarises exception counts for the datasets
with at least one discovered exception. The eight
configurations combine two mincov settings (5 and 0.1) with four precision
thresholds (0.6, 0.7, 0.8, and 0.9). Exception counts are summed over these
configurations; the final column counts configurations with a nonzero result.
They are not counts of unique exception premises after deduplication.

\begin{longtable}{@{}lrr@{}}
\caption{Classification exception counts across eight parameter configurations.}
\label{tab:classification-counts-by-dataset}\\
\toprule
Dataset & Exception rules & Configurations with exceptions \\
\midrule
\endfirsthead
\toprule
Dataset & Exception rules & Configurations with exceptions \\
\midrule
\endhead
anneal & 20 & 8 \\
mushroom & 18 & 8 \\
sonar & 8 & 6 \\
credit-g & 7 & 2 \\
vehicle & 7 & 2 \\
yeast & 6 & 2 \\
cylinder-bands & 5 & 3 \\
kdd-synthetic-control & 4 & 2 \\
diabetes & 3 & 2 \\
segment & 3 & 1 \\
ecoli & 2 & 2 \\
autos & 1 & 1 \\
bupa-liver-disorders & 1 & 1 \\
horse-colic & 1 & 1 \\
sick-euthyroid & 1 & 1 \\
\bottomrule
\end{longtable}

The table contains 87 exception-rule occurrences in 15 of the 49 datasets;
the other 34 datasets have no exceptions in the tested configurations.
Summing its exception counts reproduces the sum of the eight cells in the
total-exception table above. Anneal and Mushroom yield exceptions in all eight
configurations, whereas other datasets are more sensitive to the settings.
These counts include exception rules only; commonsense rules are excluded.

\subsubsection{Examples of classification exception rules}
\label{classification-rule-examples}
The next table presents four classification exception structures
discovered with mincov = 5 and $\theta=0.8$.  Mushroom classes p and e
denote poisonous and edible, respectively. In the count columns, $p$ and $n$ are
the covered observations belonging to the rule's target class and other classes;
$P$ and $N$ are the corresponding totals in the entire dataset. Each rule is
evaluated relative to its own conclusion, so these totals change when the
target class changes.

\begingroup
\footnotesize
\setlength{\tabcolsep}{3pt}
\begin{longtable}{@{}>{\raggedright\arraybackslash}p{0.13\linewidth}l>{\raggedright\arraybackslash}p{0.43\linewidth}rrrr@{}}
\caption{Classification CR--RR--ER examples for mincov = 5 and $\theta=0.8$.}
\label{tab:classification-rule-examples}\\
\toprule
Example & Type & Rule & $p$ & $n$ & $P$ & $N$ \\
\midrule
\endfirsthead
\toprule
Example & Type & Rule & $p$ & $n$ & $P$ & $N$ \\
\midrule
\endhead
anneal (1) & CR & IF surface-quality = \{E\} THEN class = 3 & 297 & 13 & 684 & 214 \\
anneal (1) & RR & IF carbon $\geq$ 3.50 THEN class = 3 & 74 & 2 & 684 & 214 \\
anneal (1) & ER & IF surface-quality = \{E\} AND carbon $\geq$ 3.50 THEN class = U & 2 & 0 & 40 & 858 \\
\addlinespace
anneal (2) & CR & IF carbon $\geq$ 3.50 THEN class = 3 & 74 & 2 & 684 & 214 \\
anneal (2) & RR & IF surface-quality = \{E\} THEN class = 3 & 297 & 13 & 684 & 214 \\
anneal (2) & ER & IF carbon $\geq$ 3.50 AND surface-quality = \{E\} THEN class = U & 2 & 0 & 40 & 858 \\
\addlinespace
mushroom (1) & CR & IF gill-size = \{n\} THEN class = p & 2224 & 288 & 3916 & 4208 \\
mushroom (1) & RR & IF spore-print-color = \{h\} THEN class = p & 1584 & 48 & 3916 & 4208 \\
mushroom (1) & ER & IF gill-size = \{n\} AND spore-print-color = \{h\} THEN class = e & 48 & 0 & 4208 & 3916 \\
\addlinespace
mushroom (2) & CR & IF odor = \{n\} THEN class = e & 3408 & 120 & 4208 & 3916 \\
mushroom (2) & RR & IF stalk-root = \{c\} THEN class = e & 512 & 44 & 4208 & 3916 \\
mushroom (2) & ER & IF odor = \{n\} AND stalk-root = \{c\} THEN class = p & 8 & 0 & 3916 & 4208 \\
\bottomrule
\end{longtable}
\endgroup

The two Anneal structures describe the same intersection with CR and RR
exchanged; they should not be interpreted as two distinct interactions.
Each parent rule supports class 3, while their conjunction covers two class-U
observations. Mushroom (1) reverses class p to e in an intersection containing
48 edible observations, and Mushroom (2) reverses class e to p in an
intersection containing eight poisonous observations.

All four reported ERs have no covered negatives, but two of them describe
the same subgroup and some have very small support. Perfect precision on the
discovery data does not establish generalisation or replication. In particular,
mincov controls newly covered positives during parent-rule induction and is
not a lower bound of five on the ER subgroup size.




\subsection{Regression}
\subsubsection{Discovery across datasets}
\label{regression}

\label{discovery-across-datasets-1}
Exception rules were discovered in \textbf{17 of 48} regression datasets.
The table includes every dataset with a nonzero exception
count. Rules and exceptions are counted separately; an exception count
need not represent an additional independent top-level rule.


\begingroup
\small
\setlength{\tabcolsep}{4pt}
\begin{longtable}{>{\raggedright\arraybackslash}p{\dimexpr 0.50000\linewidth-3.0000\tabcolsep\relax}>{\raggedright\arraybackslash}p{\dimexpr 0.25000\linewidth-1.5000\tabcolsep\relax}>{\raggedright\arraybackslash}p{\dimexpr 0.25000\linewidth-1.5000\tabcolsep\relax}}
\caption{Regression discovery results (datasets with exceptions).}\\
\toprule
Dataset & Rules & Exceptions \\
\midrule
\endfirsthead
\toprule
Dataset & Rules & Exceptions \\
\midrule
\endhead
bodyfat & 36 & 1 \\
breasttumor & 50 & 1 \\
cholesterol & 44 & 11 \\
compressive & 61 & 20 \\
concrete & 11 & 7 \\
friedman & 68 & 48 \\
kidney & 13 & 1 \\
laser & 97 & 3 \\
pbc & 65 & 11 \\
plastic & 188 & 24 \\
pollution & 9 & 2 \\
pwlinear & 36 & 1 \\
sensory & 91 & 3 \\
stock & 67 & 9 \\
treasury & 79 & 3 \\
triazines & 20 & 1 \\
veteran & 16 & 1 \\
\bottomrule
\end{longtable}
\endgroup
\subsubsection{Discovered rule premises}
\label{regression-rule-premises}
Each row defines a triple: CR has premise $A$, RR has premise $B$,
and ER has premise $A\land B$. This table gives one illustrated triple
for each of the 17 datasets with discovered exceptions. The aggregate
counts above may include additional exceptions for the same dataset.

\begingroup
\footnotesize
\setlength{\tabcolsep}{3pt}
\begin{longtable}{@{}p{0.04\linewidth}p{0.18\linewidth}p{0.35\linewidth}p{0.35\linewidth}@{}}
\caption{Examples of regression rule premises.}\\
\toprule
No. & Dataset & CR premise $A$ & RR premise $B$ \\
\midrule
\endfirsthead
\toprule
No. & Dataset & CR premise $A$ & RR premise $B$ \\
\midrule
\endhead
1 & bodyfat & Density $\leq$ 1.04 AND Weight $\leq$ 221.12 AND Hip $\leq$ 113.95 & Hip $\geq$ 112.15 \\
2 & breasttumor & Age $\leq$ 64.50 AND Age $\geq$ 60.50 & Age $\leq$ 61.50 \\
3 & cholesterol & Thalach $\geq$ 150.50 AND Age $\geq$ 39.50 AND Thalach $\geq$ 156.50 & Age $\geq$ 66.50 \\
4 & compressive & Age $\geq$ 73.00 & Superplasticizer $\geq$ 8.04 \\
5 & concrete & Cement $\leq$ 162.50 & Fly ash $\leq$ 155.00 \\
6 & friedman & Input1 $\leq$ 0.13 & Input4 $\leq$ 0.05 \\
7 & kidney & Age $\leq$ 45.50 AND Age $\geq$ 25.00 & Age $\leq$ 29.00 \\
8 & laser & Input1 $\geq$ 44.50 AND Input2 $\leq$ 44.50 AND Input3 $\leq$ 153.50 AND Input3 $\geq$ 8.00 AND Input4 $\leq$ 95.50 & Input2 $\leq$ 9.50 \\
9 & pbc & Z11 $\leq$ 17.50 & Z2 $\leq$ 46.50 \\
10 & plastic & Strength $\leq$ 24.35 AND Strength $\geq$ 18.15 AND Temperature $\leq$ 205.00 & Strength $\geq$ 23.10 \\
11 & pollution & NONW $\leq$ 9.15 & JANT $\geq$ 39.50 \\
12 & pwlinear & a1 $\leq$ 0.00 AND a9 $\leq$ 0.50 AND a3 $\geq$ -0.50 AND a7 $\leq$ 0.50 & a2 $\geq$ 0.50 \\
13 & sensory & Rows = \{2\} AND Interval = \{2\} AND Judges = \{4\} & Trellis = \{3\} \\
14 & stock & Company7 $\geq$ 80.19 & Company8 $\leq$ 26.06 \\
15 & treasury & Currency $\geq$ 3474.10 & CheckableDeposits $\geq$ 1148.30 \\
16 & triazines & p4\_branch $\geq$ 0.40 & p4\_flex $\leq$ 0.16 \\
17 & veteran & Age $\geq$ 58.50 AND Karnofsky $\geq$ 55.00 AND Karnofsky $\leq$ 77.50 & Months $\geq$ 13.50 \\
\bottomrule
\end{longtable}
\endgroup



\begingroup
\scriptsize
\setlength{\tabcolsep}{2pt}
\renewcommand{\arraystretch}{1.15}

\begin{longtable}{@{}c
  *{3}{>{\centering\arraybackslash}p{0.073\linewidth}}
  *{3}{>{\centering\arraybackslash}p{0.073\linewidth}}
  *{3}{>{\centering\arraybackslash}p{0.073\linewidth}}
  *{3}{>{\centering\arraybackslash}p{0.085\linewidth}}@{}}

\caption{Statistical characteristics and coverage of regression rule triples.}
\label{tab:regression_statistics}\\

\toprule
No. &
\multicolumn{3}{c}{CR} &
\multicolumn{3}{c}{RR} &
\multicolumn{3}{c}{ER} &
\multicolumn{3}{c}{Comparison $p$-values} \\
\cmidrule(lr){2-4}
\cmidrule(lr){5-7}
\cmidrule(lr){8-10}
\cmidrule(l){11-13}
&
Mean & Band & Cov. &
Mean & Band & Cov. &
Mean & Band & Cov. &
CR/RR & ER/CR & ER/RR \\
\midrule
\endfirsthead

\toprule
No. &
\multicolumn{3}{c}{CR} &
\multicolumn{3}{c}{RR} &
\multicolumn{3}{c}{ER} &
\multicolumn{3}{c}{Comparison $p$-values} \\
\cmidrule(lr){2-4}
\cmidrule(lr){5-7}
\cmidrule(lr){8-10}
\cmidrule(l){11-13}
&
Mean & Band & Cov. &
Mean & Band & Cov. &
Mean & Band & Cov. &
CR/RR & ER/CR & ER/RR \\
\midrule
\endhead

1 & 29.14 & [25.28, 33.00] & 50
  & 32.86 & [26.02, 39.71] & 13
  & 43.80 & [40.10, 47.50] & 2
  & 0.067249 & 0.018653 & 0.025641 \\

2 & 24.89 & [15.13, 34.66] & 28
  & 24.78 & [14.59, 34.97] & 233
  & 14.50 & [9.94, 19.06] & 4
  & 0.557321 & 0.016566 & 0.020288 \\

3 & 249.93 & [194.20, 305.66] & 123
  & 259.96 & [186.51, 333.41] & 26
  & 340.20 & [226.81, 453.59] & 5
  & 0.913350 & 0.013199 & 0.013378 \\

4 & 46.97 & [34.22, 59.72] & 190
  & 43.11 & [25.61, 60.62] & 400
  & 60.75 & [49.14, 72.36] & 46
  & 0.078053 & $4.25\times10^{-15}$
  & $5.77\times10^{-13}$ \\

5 & 31.69 & [26.00, 37.37] & 42
  & 33.18 & [25.99, 40.37] & 48
  & 25.16 & [20.12, 30.20] & 10
  & 0.307769 & 0.000187 & 0.000057 \\

6 & 9.52 & [5.89, 13.15] & 157
  & 8.87 & [4.94, 12.80] & 64
  & 4.88 & [2.36, 7.41] & 14
  & 0.942553 & $2.81\times10^{-6}$
  & $1.88\times10^{-5}$ \\

7 & 0.96 & [0.48, 1.44] & 28
  & 1.43 & [0.77, 2.10] & 12
  & 2.30 & [2.30, 2.30] & 2
  & 0.091689 & 0.021774 & 0.036368 \\

8 & 115.25 & [74.13, 156.36] & 166
  & 113.52 & [32.42, 194.63] & 48
  & 203.44 & [168.87, 238.02] & 9
  & 0.050321 & $2.60\times10^{-6}$
  & 0.000296 \\

9 & 2644.38 & [1944.53, 3344.22] & 16
  & 1999.96 & [911.88, 3088.05] & 156
  & 3487.25 & [2969.43, 4005.07] & 4
  & 0.050290 & 0.004396 & 0.009139 \\

10 & 13.85 & [11.36, 16.34] & 80
   & 14.39 & [11.05, 17.74] & 1319
   & 11.00 & [9.75, 12.25] & 14
   & 0.683115 & $1.67\times10^{-6}$
   & $8.26\times10^{-5}$ \\

11 & 905.15 & [855.53, 954.76] & 29
   & 934.44 & [852.57, 1016.32] & 16
   & 847.88 & [808.62, 887.15] & 4
   & 0.053553 & 0.016336 & 0.019780 \\

12 & 3.79 & [0.98, 6.59] & 26
   & 1.20 & [-4.35, 6.75] & 55
   & 7.54 & [4.74, 10.33] & 6
   & 0.099626 & 0.000556 & 0.002520 \\

13 & 15.88 & [15.20, 16.55] & 16
   & 15.28 & [14.48, 16.09] & 144
   & 16.75 & [16.19, 17.31] & 4
   & 0.077285 & 0.007555 & 0.002585 \\

14 & 48.17 & [44.51, 51.84] & 137
   & 46.20 & [39.36, 53.04] & 691
   & 53.45 & [50.12, 56.78] & 33
   & 0.117571 & $1.34\times10^{-12}$
   & $3.35\times10^{-10}$ \\

15 & 5.46 & [5.17, 5.76] & 250
   & 5.31 & [4.70, 5.92] & 35
   & 5.94 & [5.83, 6.05] & 4
   & 0.054602 & 0.002003 & 0.048704 \\

16 & 0.58 & [0.43, 0.73] & 37
   & 0.61 & [0.40, 0.82] & 67
   & 0.10 & [0.10, 0.10] & 2
   & 0.235890 & 0.020417 & 0.017455 \\

17 & 105.46 & [21.27, 189.65] & 35
   & 112.90 & [17.49, 208.32] & 21
   & 208.75 & [154.60, 262.90] & 4
   & 0.483457 & 0.005530 & 0.035257 \\

\bottomrule
\end{longtable}
\endgroup

Table~\ref{tab:regression_statistics} summarizes the statistical characteristics of the 17 illustrated regression rule triples, using the same numbering as the preceding premise table. For each triple, it reports the mean target value, the associated variability band, and the number of covered observations for CR, RR, and ER, together with the three pairwise comparison $p$-values. The bands describe within-group variability rather than confidence intervals for the estimated means.

Across all 17 examples, the CR and RR comparisons yield $p$-values above 0.05, whereas both ER comparisons yield values below 0.05. Thus, under the applied acceptance criteria, the parent groups are statistically compatible while their intersection exhibits a significant deviation from both. The direction of this deviation varies across datasets: the ER mean exceeds both parent means in nine examples and falls below both in the remaining eight. This indicates that the discovered exceptions capture both unusually high and unusually low target values, rather than deviations in a single direction.

The coverage statistics further illustrate the localized nature of these patterns. ER coverage ranges from 2 to 46 observations and is consistently smaller than the coverage of either parent rule. Several examples involve very small intersections, particularly bodyfat, kidney, and triazines, each with only two ER observations. Consequently, although the reported comparisons satisfy the discovery criteria, patterns with limited support require particular caution and independent validation. The results should therefore be interpreted as evidence of statistically distinguishable subgroups within the discovery datasets, rather than as confirmation of their generalizability or causal significance.

\subsection{Survival analysis}
\subsubsection{Discovery across datasets}
\label{survival-analysis}

\label{discovery-across-datasets-2}
Exception rules were discovered in \textbf{18 of 25} survival datasets.
The table below includes all 25 datasets, including those
where no exceptions were found.


\begingroup
\small
\setlength{\tabcolsep}{4pt}
\begin{longtable}{>{\raggedright\arraybackslash}p{\dimexpr 0.50000\linewidth-3.0000\tabcolsep\relax}>{\raggedright\arraybackslash}p{\dimexpr 0.25000\linewidth-1.5000\tabcolsep\relax}>{\raggedright\arraybackslash}p{\dimexpr 0.25000\linewidth-1.5000\tabcolsep\relax}}
\caption{Survival discovery results.}\\
\toprule
Dataset & Rules & Exceptions \\
\midrule
\endfirsthead
\toprule
Dataset & Rules & Exceptions \\
\midrule
\endhead
actg320 & 4 & 1 \\
BHS & 4 & 0 \\
bmt-ch & 2 & 0 \\
cancer & 11 & 0 \\
csb-ch07-halibut & 3 & 1 \\
csb-ch12-nursing & 15 & 3 \\
csb-ch15-lung & 2 & 1 \\
follic & 2 & 1 \\
GBSG2 & 2 & 1 \\
hd & 2 & 0 \\
lac\_selected & 3 & 2 \\
lung & 2 & 1 \\
Melanoma & 9 & 1 \\
mgus & 5 & 1 \\
pbc & 3 & 1 \\
przeszczepy & 2 & 0 \\
Rossi & 4 & 1 \\
std & 6 & 3 \\
TRACE & 2 & 0 \\
uis & 2 & 1 \\
unemployed & 9 & 1 \\
wcgs & 3 & 3 \\
whas1 & 3 & 3 \\
whas500 & 3 & 2 \\
zinc & 3 & 0 \\
\bottomrule
\end{longtable}
\endgroup
\subsubsection{Discovered rule premises}
\label{survival-rule-premises}
Each row defines a triple: CR has premise $A$, RR has premise $B$,
and ER has premise $A\land B$. This table lists 28 discovered triples,
complementing the dataset-level counts above.

\begingroup
\footnotesize
\setlength{\tabcolsep}{3pt}
\begin{longtable}{@{}p{0.04\linewidth}p{0.18\linewidth}p{0.35\linewidth}p{0.35\linewidth}@{}}
\caption{Discovered survival rule premises.}\\
\toprule
No. & Dataset & CR premise $A$ & RR premise $B$ \\
\midrule
\endfirsthead
\toprule
No. & Dataset & CR premise $A$ & RR premise $B$ \\
\midrule
\endhead
1 & actg320 & cd4 $\leq$ 18.75 & karnof $\leq$ 85.00 \\
2 & csb-ch07-halibut & HANDTIME $\leq$ 15.50 & LOGCAT $\geq$ 4.59 \\
3 & csb-ch12-nursing & gender = \{0\} AND age $\geq$ 73.50 & hlstat $\leq$ 3.50 AND age $\geq$ 67.50 \\
4 & csb-ch12-nursing & gender = \{0\} AND age $\geq$ 71.50 & hlstat $\leq$ 3.50 AND age $\geq$ 65.50 \\
5 & csb-ch12-nursing & gender = \{0\} AND age $\geq$ 67.50 & hlstat $\leq$ 3.50 \\
6 & csb-ch15-lung & Stage4 = \{0\} & Cell\_Type $\leq$ 1.50 \\
7 & follic & age $\geq$ 66.95 & hgb $\leq$ 117.50 AND age $\geq$ 26.05 \\
8 & GBSG2 & pnodes $\geq$ 4.50 & progrec $\leq$ 2.50 \\
9 & lac\_selected & INHA\_a $\leq$ 11.33 AND UBC $\leq$ 12.51 & SLC1A6 $\leq$ 9.25 AND CRK $\leq$ 8.57 \\
10 & lac\_selected & INHA\_a $\leq$ 11.33 AND FADD $\leq$ 10.24 AND ATP2B1 $\leq$ 9.92 & GCNT1\_a $\leq$ 7.50 \\
11 & lung & Stage4 = \{0\} & Cell\_Type $\leq$ 1.50 \\
12 & Melanoma & thick $\leq$ 3.38 AND age $\geq$ 14.50 & age $\leq$ 55.50 AND thick $\leq$ 12.16 \\
13 & mgus & age $\geq$ 62.50 & hgb $\leq$ 12.85 \\
14 & pbc & bili $\geq$ 1.05 & albumin $\leq$ 3.55 \\
15 & Rossi & age $\leq$ 19.50 & prio $\geq$ 8.50 AND emp6 = \{no\} \\
16 & std & age $\geq$ 17.50 AND vagina $\leq$ 0.50 AND npartner $\leq$ 3.50 AND iinfct $\geq$ 1.50 & race = \{W\} \\
17 & std & age $\geq$ 17.50 AND vagina $\leq$ 0.50 AND age $\leq$ 42.00 AND iinfct $\geq$ 1.50 & race = \{W\} \\
18 & std & age $\leq$ 17.50 & dchexam $\leq$ 0.50 \\
19 & uis & FRAC $\leq$ 0.94 & NDT $\geq$ 5.50 AND FRAC $\leq$ 1.56 \\
20 & unemployed & age = \{middle\} & gender = \{male\} AND birth $\leq$ 205.50 \\
21 & wcgs & chol0 $\leq$ 271.50 & dbp0 $\leq$ 84.50 \\
22 & wcgs & chol0 $\geq$ 264.50 & sbp0 $\geq$ 147.00 \\
23 & wcgs & age0 $\leq$ 47.50 & chol0 $\leq$ 236.50 AND sbp0 $\leq$ 186.00 \\
24 & whas1 & SHO $\leq$ 0.50 & CPK $\geq$ 701.50 \\
25 & whas1 & AGE $\leq$ 69.50 & CHF $\leq$ 0.50 \\
26 & whas1 & AGE $\geq$ 66.50 & CPK $\geq$ 2160.00 AND AGE $\geq$ 36.50 \\
27 & whas500 & chf $\leq$ 0.50 & hr $\leq$ 71.50 \\
28 & whas500 & age $\geq$ 64.50 & bmi $\leq$ 25.40 \\
\bottomrule
\end{longtable}
\endgroup






% Editable supplement chapter.
\section{Computational efficiency}
\label{supplement-computational-efficiency}

This section collects induction times for classification, regression, and
survival analysis. For each task, the comparison includes the exception-enabled
learner, its version without exception search, RuleKit, and a task-specific
decision tree. Classification also includes a parameter-sensitivity study. All durations are
reported in seconds. Measurement scope is specified for each task;
the three task-specific experiments do not form a common-platform speed benchmark.

\subsection{Classification}
\subsubsection{Parameter sensitivity}
The table below summarises median induction times across parameter settings.
\label{execution-times}

\begingroup
\small
\setlength{\tabcolsep}{4pt}
\begin{longtable}{>{\raggedright\arraybackslash}p{\dimexpr 0.20000\linewidth-2.0000\tabcolsep\relax}>{\raggedright\arraybackslash}p{\dimexpr 0.20000\linewidth-2.0000\tabcolsep\relax}>{\raggedright\arraybackslash}p{\dimexpr 0.20000\linewidth-2.0000\tabcolsep\relax}>{\raggedright\arraybackslash}p{\dimexpr 0.20000\linewidth-2.0000\tabcolsep\relax}>{\raggedright\arraybackslash}p{\dimexpr 0.20000\linewidth-2.0000\tabcolsep\relax}}
\caption{Median classification exception-rule generation times (seconds).}\\
\toprule
mincov & theta = 0.6 & theta = 0.7 & theta = 0.8 & theta = 0.9 \\
\midrule
\endfirsthead
\toprule
mincov & theta = 0.6 & theta = 0.7 & theta = 0.8 & theta = 0.9 \\
\midrule
\endhead
5 & 66 & 86 & 23 & 4 \\
0.1 & 22 & 9 & 3 & 2 \\
\bottomrule
\end{longtable}
\endgroup
The absolute-coverage setting is slower at every tested threshold. Within a
setting the runtime is not monotonic in theta: with mincov 5, the median rises
from 66 to 86 seconds before falling. The number of candidate refinements and
early stopping affect runtime in addition to dataset dimensions.





\subsubsection{Comparison with other induction methods}
Four methods are compared on the same 49 datasets, using one fitting-time
observation per method and dataset (196 observations). RulesClassifier
and ExceptionRulesClassifier denote the learners without and with
exception search, respectively.

\begin{table}[htbp]
\centering
\small
\setlength{\tabcolsep}{4pt}
\caption{Classification fitting times in seconds, aggregated from saved comparison records (49 datasets per method).}
\label{tab:classification-fitting-times}
\begin{tabularx}{\linewidth}{@{}Xrrrrrr@{}}
\toprule
Method & Mean & Median & Q1 & Q3 & Min & Max \\
\midrule
DecisionTreeClassifier & 0.025 & 0.014 & 0.008 & 0.021 & 0.005 & 0.170 \\
RuleKit & 0.437 & 0.291 & 0.127 & 0.459 & 0.013 & 2.602 \\
RulesClassifier & 5.661 & 0.766 & 0.267 & 2.589 & 0.010 & 78.187 \\
ExceptionRulesClassifier & 43.338 & 3.822 & 1.201 & 16.734 & 0.010 & 1009.613 \\
\bottomrule
\end{tabularx}
\end{table}

The four prediction variants of the exception-enabled model share one fit;
its fitting time is counted once per dataset.

The comparison uses a 70/30 train/test split with \texttt{random\_state=42},
minimum new coverage of five, C2 induction quality, and a five-condition growth
limit; the classification tree has maximum depth five. Elapsed wall-clock
time is measured with \texttt{time.time()} around the wrapper's \texttt{fit} call.
For the tree this includes preprocessing; for RuleKit it includes estimator
construction, fitting, and conversion to package rules. Data loading, splitting,
subsequent prediction, and result export are outside the interval. CPU times
are not included in the table.

Means and quantiles weight datasets equally. Q1 and Q3 use inclusive linear
interpolation, as in the survival summary. The median rises from 0.766 seconds
without exception search to 3.822 seconds with it. RuleKit and the classification
tree are faster in these saved runs. These measurements concern
a training-partition experiment, separate from the complete-dataset parameter
study above.


\subsection{Regression}
The table below summarises fitting times for the regression comparison.
\label{training-times}
The timed methods are the regression learner with and without exception
search, RuleKit, and DecisionTreeRegressor. Rule growth and tree depth are
limited to five. The comparison uses the training partition
of a 70/30 split with \texttt{random\_state=42}; these fitting times are separate
from the complete-dataset discovery counts above.

\begingroup
\small
\setlength{\tabcolsep}{4pt}
\begin{longtable}{>{\raggedright\arraybackslash}p{\dimexpr 0.28000\linewidth-3.9200\tabcolsep\relax}>{\raggedright\arraybackslash}p{\dimexpr 0.12000\linewidth-1.6800\tabcolsep\relax}>{\raggedright\arraybackslash}p{\dimexpr 0.12000\linewidth-1.6800\tabcolsep\relax}>{\raggedright\arraybackslash}p{\dimexpr 0.12000\linewidth-1.6800\tabcolsep\relax}>{\raggedright\arraybackslash}p{\dimexpr 0.12000\linewidth-1.6800\tabcolsep\relax}>{\raggedright\arraybackslash}p{\dimexpr 0.12000\linewidth-1.6800\tabcolsep\relax}>{\raggedright\arraybackslash}p{\dimexpr 0.12000\linewidth-1.6800\tabcolsep\relax}}
\caption{Regression training times in seconds.}\\
\toprule
Method & Mean & Median & Q1 & Q3 & Min & Max \\
\midrule
\endfirsthead
\toprule
Method & Mean & Median & Q1 & Q3 & Min & Max \\
\midrule
\endhead
Decision\allowbreak Tree\allowbreak Regressor & 0.009 & 0.008 & 0.006 & 0.010 & 0.004 & 0.017 \\
RuleKit & 0.696 & 0.182 & 0.062 & 1.298 & 0.010 & 2.951 \\
RulesRegressor & 15.414 & 1.148 & 0.366 & 9.028 & 0.028 & 166.118 \\
ExceptionRulesRegressor & 77.055 & 4.863 & 1.345 & 41.804 & 0.080 & 975.515 \\
\bottomrule
\end{longtable}
\endgroup
Exception search raises the reported median from 1.148 to 4.863 seconds and
the maximum from 166.118 to 975.515 seconds. The tree and RuleKit are faster.
This reflects the added candidate search and statistical checks, together with
implementation differences; it should not be read as an isolated measurement
of algorithmic complexity. The large difference between mean and median shows
substantial variability across datasets.

\subsection{Survival analysis}
\label{execution-times-1}
Table~\ref{tab:survival-fitting-times} summarises 100 fitting-time records:
one observation per method and dataset for four methods on 25 datasets.

Models were fitted on the training partition of a 70/30 split
with \texttt{random\_state=42}. Rule growth and survival-tree depth are limited
to five. These times are separate from complete-dataset discovery counts.

\begin{table}[htbp]
\centering
\small
\setlength{\tabcolsep}{4pt}
\caption{Survival fitting times in seconds, aggregated from saved experiment records (25 datasets per method).}
\label{tab:survival-fitting-times}
\begin{tabularx}{\linewidth}{@{}Xrrrrrr@{}}
\toprule
Method & Mean & Median & Q1 & Q3 & Min & Max \\
\midrule
SurvivalTree & 0.025 & 0.015 & 0.012 & 0.022 & 0.007 & 0.171 \\
RuleKit & 0.687 & 0.445 & 0.358 & 0.793 & 0.181 & 2.029 \\
RulesSurvival & 69.754 & 26.366 & 12.175 & 53.005 & 2.517 & 399.441 \\
ExceptionRulesSurvival & 274.255 & 97.846 & 42.313 & 193.187 & 4.240 & 1933.616 \\
\bottomrule
\end{tabularx}
\end{table}

The comparison script measures elapsed wall-clock time using \texttt{time.time()}
immediately before and after the model wrapper's \texttt{fit} call. This includes
work inside that wrapper: preprocessing for the survival tree, and model
construction, fitting, and conversion to package rules for RuleKit. Dataset
loading, the train/test split, subsequent predictive evaluation, and result
export lie outside this interval. The separately saved \texttt{process\_time} and
\texttt{thread\_time} CPU measurements are not used in the table.

The summaries weight every dataset equally and use linear-interpolated 25th
and 75th percentiles (the inclusive convention). They describe variation across
datasets rather than repeated-run uncertainty.

Exception search increases the median from 26.366 to 97.846 seconds. The ratio
of the unrounded medians is approximately 3.71; this is a ratio of aggregate
medians, not a median of per-dataset slowdowns. The largest exception-enabled
time is 1933.616 seconds (about 32.2 minutes) on TRACE, compared with 399.441
seconds without exceptions on that dataset. SurvivalTree and RuleKit are
faster in this saved comparison. The difference between means and medians
indicates an uneven distribution of computational costs across datasets.


\subsection{Dataset size and induction time}
\label{runtime-by-size}

The dataset-level analysis combines fitting times with the dimensions of the
49 classification, 48 regression, and 25 survival benchmark datasets.

M is the full dataset's number of rows (observations), while N denotes the
number of descriptive columns before encoding.
The total number of columns, including outcomes, is denoted by
$N_{\mathrm{all}}$.
For classification and regression, N excludes class. For survival, N excludes
survival\_time, survival\_status, and the optional group column, in accordance
with the induction pipeline. The timing scripts use training partitions from
a 70/30 split, so fitted row counts are M minus $\lceil 0.3M \rceil$.
M and $N_{\mathrm{all}}$ describe the input datasets before encoding.

Tree denotes
the task-specific decision tree, Rules the learner without exception search,
and Exceptions the exception-enabled learner. Results are sorted by M, then N.
\begingroup
\footnotesize
\setlength{\tabcolsep}{3pt}
\begin{longtable}{@{}llrrr@{}}
\caption{Spearman rank correlations between fitting time and dataset dimensions.}\\
\toprule
Task & Method & Datasets & $\rho(t,M)$ & $\rho(t,N)$ \\
\midrule
\endfirsthead
\toprule
Task & Method & Datasets & $\rho(t,M)$ & $\rho(t,N)$ \\
\midrule
\endhead
Classification & Tree & 49 & 0.530 & 0.869 \\
Classification & RuleKit & 49 & 0.708 & 0.584 \\
Classification & Rules & 49 & 0.683 & 0.451 \\
Classification & Exceptions & 49 & 0.530 & 0.389 \\
Regression & Tree & 48 & 0.346 & 0.787 \\
Regression & RuleKit & 48 & 0.923 & 0.319 \\
Regression & Rules & 48 & 0.874 & 0.325 \\
Regression & Exceptions & 48 & 0.859 & 0.336 \\
Survival analysis & Tree & 25 & 0.207 & 0.538 \\
Survival analysis & RuleKit & 25 & 0.625 & 0.037 \\
Survival analysis & Rules & 25 & 0.291 & 0.174 \\
Survival analysis & Exceptions & 25 & 0.351 & 0.248 \\
\bottomrule
\end{longtable}
\endgroup

Spearman correlations use average ranks for ties and quantify separate
associations with M and N within each task and method. They do not adjust for
the other dimension, attribute cardinality, rule-search complexity, or censoring.
These are descriptive comparisons across different datasets, not a controlled
scaling experiment or an estimate of algorithmic complexity. The per-dataset
tables below retain the individual observations behind these summaries.

For the exception-enabled learner, the association with row count is strongest
in regression (rho = 0.859), followed by classification (0.530) and survival
(0.351). The respective associations with attributes count are 0.336, 0.389,
and 0.248. 


\subsubsection{Classification datasets}
\begingroup
\footnotesize
\setlength{\tabcolsep}{3pt}
\begin{longtable}{@{}p{0.25\linewidth}rrrrrr@{}}
\caption{Classification: dataset dimensions and saved fitting times (seconds).}\\
\toprule
Dataset & M & N & Tree & RuleKit & Rules & Exceptions \\
\midrule
\endfirsthead
\toprule
Dataset & M & N & Tree & RuleKit & Rules & Exceptions \\
\midrule
\endhead
\nolinkurl{contact-lenses} & 24 & 4 & 0.006 & 0.013 & 0.010 & 0.010 \\
\nolinkurl{labor} & 57 & 16 & 0.011 & 0.194 & 0.027 & 0.322 \\
\nolinkurl{zoo} & 101 & 16 & 0.013 & 0.056 & 0.051 & 0.051 \\
\nolinkurl{echocardiogram} & 131 & 11 & 0.008 & 0.056 & 0.095 & 0.748 \\
\nolinkurl{hayes-roth} & 132 & 4 & 0.006 & 0.038 & 0.028 & 0.063 \\
\nolinkurl{lymph} & 148 & 18 & 0.014 & 0.071 & 0.122 & 0.340 \\
\nolinkurl{iris} & 150 & 4 & 0.005 & 0.046 & 0.064 & 0.827 \\
\nolinkurl{hepatitis} & 155 & 19 & 0.014 & 0.228 & 0.164 & 1.283 \\
\nolinkurl{wine} & 178 & 13 & 0.007 & 0.060 & 0.311 & 3.264 \\
\nolinkurl{flag} & 194 & 28 & 0.017 & 0.318 & 0.417 & 2.359 \\
\nolinkurl{autos} & 205 & 25 & 0.017 & 0.380 & 1.434 & 12.123 \\
\nolinkurl{sonar} & 208 & 60 & 0.020 & 0.172 & 5.358 & 50.191 \\
\nolinkurl{glass} & 214 & 9 & 0.007 & 0.213 & 0.941 & 6.913 \\
\nolinkurl{audiology} & 226 & 69 & 0.039 & 0.519 & 0.510 & 1.211 \\
\nolinkurl{prnn-synth} & 250 & 2 & 0.005 & 0.061 & 0.153 & 1.214 \\
\nolinkurl{heart-statlog} & 270 & 13 & 0.007 & 0.233 & 0.698 & 4.662 \\
\nolinkurl{breast-cancer} & 286 & 9 & 0.013 & 0.094 & 0.187 & 0.345 \\
\nolinkurl{hungarian-heart-disease} & 294 & 13 & 0.009 & 0.246 & 0.435 & 2.576 \\
\nolinkurl{cleveland} & 303 & 13 & 0.010 & 0.286 & 1.198 & 5.822 \\
\nolinkurl{heart-c} & 303 & 13 & 0.013 & 0.242 & 0.749 & 3.936 \\
\nolinkurl{ecoli} & 336 & 7 & 0.006 & 0.223 & 0.708 & 13.795 \\
\nolinkurl{primary-tumor} & 339 & 17 & 0.015 & 0.149 & 0.525 & 0.944 \\
\nolinkurl{bupa-liver-disorders} & 345 & 6 & 0.006 & 0.294 & 1.236 & 8.508 \\
\nolinkurl{ionosphere} & 351 & 34 & 0.015 & 0.134 & 2.589 & 15.161 \\
\nolinkurl{horse-colic} & 368 & 22 & 0.018 & 0.459 & 0.863 & 3.822 \\
\nolinkurl{auto-mpg} & 398 & 7 & 0.007 & 0.443 & 1.164 & 11.092 \\
\nolinkurl{vote} & 435 & 16 & 0.015 & 0.078 & 0.055 & 0.239 \\
\nolinkurl{cylinder-bands} & 540 & 35 & 0.026 & 0.435 & 4.908 & 29.324 \\
\nolinkurl{kdd-synthetic-control} & 600 & 60 & 0.030 & 0.865 & 40.641 & 394.847 \\
\nolinkurl{balance-scale} & 625 & 4 & 0.006 & 0.444 & 0.349 & 1.061 \\
\nolinkurl{soybean} & 683 & 35 & 0.036 & 0.453 & 0.692 & 2.888 \\
\nolinkurl{credit-a} & 690 & 15 & 0.017 & 0.388 & 2.338 & 16.734 \\
\nolinkurl{breast-w} & 699 & 9 & 0.007 & 0.382 & 0.267 & 3.424 \\
\nolinkurl{diabetes} & 768 & 8 & 0.008 & 0.476 & 7.397 & 32.072 \\
\nolinkurl{vehicle} & 846 & 18 & 0.012 & 0.390 & 15.178 & 116.268 \\
\nolinkurl{anneal} & 898 & 38 & 0.040 & 2.035 & 0.842 & 2.347 \\
\nolinkurl{tic-tac-toe} & 958 & 9 & 0.014 & 0.127 & 0.518 & 2.091 \\
\nolinkurl{mammographic-masses} & 961 & 5 & 0.008 & 0.101 & 0.766 & 6.005 \\
\nolinkurl{credit-g} & 1000 & 20 & 0.025 & 0.570 & 7.129 & 37.547 \\
\nolinkurl{yeast} & 1484 & 8 & 0.009 & 0.291 & 16.793 & 93.715 \\
\nolinkurl{car} & 1728 & 6 & 0.016 & 0.326 & 0.862 & 1.201 \\
\nolinkurl{titanic} & 2201 & 3 & 0.012 & 0.074 & 0.165 & 0.317 \\
\nolinkurl{segment} & 2310 & 19 & 0.021 & 0.758 & 78.187 & 1009.613 \\
\nolinkurl{seismic-bumps} & 2584 & 18 & 0.025 & 0.959 & 50.086 & 89.303 \\
\nolinkurl{splice} & 3190 & 60 & 0.170 & 2.602 & 7.116 & 42.044 \\
\nolinkurl{kr-vs-kp} & 3196 & 36 & 0.096 & 1.020 & 2.051 & 4.535 \\
\nolinkurl{hypothyroid} & 3772 & 29 & 0.081 & 1.222 & 5.486 & 36.884 \\
\nolinkurl{sick-euthyroid} & 3772 & 29 & 0.082 & 0.936 & 13.754 & 47.012 \\
\nolinkurl{mushroom} & 8124 & 22 & 0.139 & 1.266 & 1.782 & 2.505 \\
\bottomrule
\end{longtable}
\endgroup

\subsubsection{Regression datasets}
\begingroup
\footnotesize
\setlength{\tabcolsep}{3pt}
\begin{longtable}{@{}p{0.25\linewidth}rrrrrr@{}}
\caption{Regression: dataset dimensions and saved fitting times (seconds).}\\
\toprule
Dataset & M & N & Tree & RuleKit & Rules & Exceptions \\
\midrule
\endfirsthead
\toprule
Dataset & M & N & Tree & RuleKit & Rules & Exceptions \\
\midrule
\endhead
\nolinkurl{longley} & 16 & 6 & 0.006 & 0.010 & 0.028 & 0.080 \\
\nolinkurl{gascons} & 27 & 4 & 0.004 & 0.014 & 0.063 & 0.223 \\
\nolinkurl{bolts} & 40 & 7 & 0.005 & 0.023 & 0.081 & 0.289 \\
\nolinkurl{diabetes} & 43 & 2 & 0.004 & 0.012 & 0.058 & 0.199 \\
\nolinkurl{elusage} & 55 & 2 & 0.006 & 0.015 & 0.061 & 0.223 \\
\nolinkurl{pollution} & 60 & 15 & 0.007 & 0.056 & 0.728 & 3.155 \\
\nolinkurl{mbagrade} & 61 & 2 & 0.005 & 0.016 & 0.041 & 0.151 \\
\nolinkurl{pyrim} & 74 & 27 & 0.010 & 0.062 & 0.149 & 0.513 \\
\nolinkurl{kidney} & 76 & 5 & 0.006 & 0.057 & 0.157 & 0.683 \\
\nolinkurl{auto93} & 93 & 22 & 0.015 & 0.153 & 0.942 & 4.302 \\
\nolinkurl{concrete} & 103 & 9 & 0.008 & 0.060 & 1.035 & 2.875 \\
\nolinkurl{cloud} & 108 & 6 & 0.009 & 0.032 & 0.425 & 1.990 \\
\nolinkurl{fruitfly} & 125 & 4 & 0.005 & 0.066 & 0.240 & 1.016 \\
\nolinkurl{echomonths} & 130 & 9 & 0.008 & 0.077 & 0.937 & 4.296 \\
\nolinkurl{veteran} & 137 & 7 & 0.008 & 0.062 & 0.270 & 1.072 \\
\nolinkurl{fishcatch} & 158 & 7 & 0.006 & 0.069 & 1.261 & 6.816 \\
\nolinkurl{auto-price} & 159 & 15 & 0.009 & 0.140 & 1.292 & 5.425 \\
\nolinkurl{servo} & 167 & 4 & 0.008 & 0.040 & 0.072 & 0.182 \\
\nolinkurl{triazines} & 186 & 60 & 0.017 & 0.191 & 0.795 & 3.370 \\
\nolinkurl{lowbwt} & 189 & 9 & 0.008 & 0.065 & 0.392 & 1.436 \\
\nolinkurl{pharynx} & 195 & 11 & 0.012 & 0.174 & 0.607 & 2.534 \\
\nolinkurl{pwlinear} & 200 & 10 & 0.006 & 0.141 & 0.288 & 0.854 \\
\nolinkurl{machine} & 209 & 6 & 0.005 & 0.066 & 0.644 & 2.317 \\
\nolinkurl{cpu} & 209 & 7 & 0.011 & 0.066 & 0.575 & 2.468 \\
\nolinkurl{bodyfat} & 252 & 14 & 0.008 & 0.684 & 7.393 & 40.962 \\
\nolinkurl{breasttumor} & 286 & 9 & 0.013 & 0.343 & 0.723 & 3.062 \\
\nolinkurl{cholesterol} & 303 & 13 & 0.011 & 0.246 & 2.203 & 10.733 \\
\nolinkurl{ozone} & 330 & 8 & 0.006 & 0.665 & 3.629 & 15.894 \\
\nolinkurl{baseball} & 337 & 16 & 0.009 & 0.653 & 8.992 & 39.121 \\
\nolinkurl{dee} & 365 & 6 & 0.007 & 0.551 & 9.139 & 44.332 \\
\nolinkurl{auto-mpg} & 398 & 7 & 0.008 & 1.815 & 2.463 & 11.656 \\
\nolinkurl{pbc} & 418 & 18 & 0.012 & 0.922 & 14.635 & 84.217 \\
\nolinkurl{ele-1} & 495 & 2 & 0.005 & 0.124 & 1.596 & 7.330 \\
\nolinkurl{housing} & 506 & 13 & 0.010 & 1.638 & 16.849 & 92.095 \\
\nolinkurl{forestFires} & 517 & 12 & 0.008 & 1.166 & 7.325 & 32.995 \\
\nolinkurl{meta} & 528 & 21 & 0.014 & 0.676 & 2.074 & 7.990 \\
\nolinkurl{sensory} & 576 & 11 & 0.012 & 2.624 & 0.991 & 2.961 \\
\nolinkurl{strike} & 625 & 6 & 0.007 & 1.289 & 4.221 & 19.598 \\
\nolinkurl{stock} & 950 & 9 & 0.010 & 2.578 & 32.920 & 159.877 \\
\nolinkurl{laser} & 993 & 4 & 0.006 & 1.340 & 15.474 & 72.737 \\
\nolinkurl{compressive} & 1030 & 8 & 0.008 & 1.325 & 38.464 & 117.185 \\
\nolinkurl{mortgage} & 1049 & 15 & 0.013 & 1.353 & 125.302 & 782.804 \\
\nolinkurl{treasury} & 1049 & 15 & 0.015 & 2.951 & 166.118 & 975.515 \\
\nolinkurl{ele-2} & 1056 & 4 & 0.005 & 1.002 & 1.851 & 7.481 \\
\nolinkurl{friedman} & 1200 & 5 & 0.008 & 2.776 & 160.875 & 621.778 \\
\nolinkurl{wizmir} & 1461 & 9 & 0.010 & 1.437 & 41.206 & 190.654 \\
\nolinkurl{wankara} & 1609 & 9 & 0.010 & 2.164 & 51.580 & 238.800 \\
\nolinkurl{plastic} & 1650 & 2 & 0.007 & 1.415 & 12.726 & 72.409 \\
\bottomrule
\end{longtable}
\endgroup

\subsubsection{Survival analysis datasets}
\begingroup
\footnotesize
\setlength{\tabcolsep}{3pt}
\begin{longtable}{@{}p{0.25\linewidth}rrrrrr@{}}
\caption{Survival analysis: dataset dimensions and saved fitting times (seconds).}\\
\toprule
Dataset & M & N & Tree & RuleKit & Rules & Exceptions \\
\midrule
\endfirsthead
\toprule
Dataset & M & N & Tree & RuleKit & Rules & Exceptions \\
\midrule
\endhead
\nolinkurl{BHS} & 75 & 12 & 0.007 & 0.354 & 4.480 & 17.944 \\
\nolinkurl{lac_selected} & 86 & 113 & 0.038 & 0.793 & 247.041 & 598.523 \\
\nolinkurl{bmt-ch} & 187 & 35 & 0.022 & 0.254 & 27.275 & 122.337 \\
\nolinkurl{przeszczepy} & 187 & 37 & 0.023 & 0.417 & 27.498 & 123.616 \\
\nolinkurl{Melanoma} & 205 & 7 & 0.012 & 0.181 & 11.760 & 52.019 \\
\nolinkurl{cancer} & 228 & 7 & 0.008 & 0.267 & 12.175 & 57.252 \\
\nolinkurl{mgus} & 241 & 9 & 0.013 & 0.297 & 21.031 & 35.971 \\
\nolinkurl{csb-ch07-halibut} & 294 & 5 & 0.009 & 0.445 & 7.073 & 22.098 \\
\nolinkurl{unemployed} & 300 & 6 & 0.013 & 0.362 & 13.215 & 41.983 \\
\nolinkurl{pbc} & 418 & 17 & 0.022 & 0.461 & 89.978 & 177.589 \\
\nolinkurl{zinc} & 431 & 55 & 0.171 & 0.715 & 26.366 & 131.657 \\
\nolinkurl{Rossi} & 432 & 60 & 0.048 & 0.709 & 9.066 & 42.313 \\
\nolinkurl{whas1} & 481 & 7 & 0.010 & 0.359 & 23.501 & 86.742 \\
\nolinkurl{whas500} & 500 & 13 & 0.015 & 0.406 & 53.005 & 193.187 \\
\nolinkurl{follic} & 541 & 4 & 0.013 & 0.398 & 38.965 & 78.592 \\
\nolinkurl{uis} & 575 & 13 & 0.012 & 0.238 & 28.388 & 130.942 \\
\nolinkurl{GBSG2} & 686 & 8 & 0.022 & 0.945 & 218.451 & 654.843 \\
\nolinkurl{hd} & 865 & 6 & 0.015 & 0.358 & 15.035 & 76.288 \\
\nolinkurl{std} & 877 & 21 & 0.019 & 0.616 & 19.805 & 97.846 \\
\nolinkurl{lung} & 1032 & 7 & 0.013 & 1.901 & 2.517 & 6.097 \\
\nolinkurl{csb-ch15-lung} & 1032 & 7 & 0.012 & 0.475 & 2.631 & 4.240 \\
\nolinkurl{actg320} & 1151 & 11 & 0.016 & 1.370 & 145.388 & 1180.621 \\
\nolinkurl{csb-ch12-nursing} & 1601 & 5 & 0.014 & 0.798 & 39.658 & 201.104 \\
\nolinkurl{TRACE} & 1878 & 6 & 0.046 & 2.029 & 399.441 & 1933.616 \\
\nolinkurl{wcgs} & 3154 & 10 & 0.041 & 2.020 & 260.117 & 788.943 \\
\bottomrule
\end{longtable}
\endgroup


\section{Discussion}
\label{outstanding-evidence}
The results demonstrate the discovery of explicit CR--RR--ER structures in
classification, regression, and survival analysis. The common representation
allows each exceptional subgroup to be inspected together with both parent
rules. Classification counts vary with coverage and precision requirements,
while the regression and survival examples quantify the outcome differences
associated with the intersections.

The evaluation concerns exploratory knowledge discovery. Rule counts measure
occurrences satisfying the specified criteria; the subgroup contrasts describe
the discovery data. Statistical comparisons are used as unadjusted search
criteria. In survival analysis the compared parent and intersection groups
overlap, and non-rejection of a parent comparison is not an equivalence test.
The reported patterns are therefore interpreted as empirical associations
and hypotheses for domain investigation.

Exception search increases the recorded fitting cost relative to the base
learners. Per-dataset timings show substantial variation, and the association
with dataset size differs between tasks. The computational results describe
the recorded runs within each experiment; the dimension correlations are
descriptive associations across datasets.

\clearpage
% Editable supplement chapter.


\begin{thebibliography}{99}


\bibitem{MachaRegression2026}
D.~Macha, {\L}.~Wr{\'o}bel, and M.~Sikora,
``Regression Exception Rules,''
in \emph{Proceedings of the 21st Conference on Computer Science and Intelligence Systems (FedCSIS)},
pp.~105--113, 2026.

\bibitem{MachaSurvival2026}
D.~Macha, {\L}.~Wr{\'o}bel, and M.~Sikora,
``Survival Exception Rules,''
in \emph{Procedia Computer Science},
30th International Conference on Knowledge-Based and Intelligent Information \& Engineering Systems (KES 2026), 2026.

\bibitem{Janssen2011}
F. Janssen, J. Fürnkranz,
\newblock  Heuristic rule-based regression via dynamic reduction to classification, 
\newblock in: IJCAI’11, AAAI Press, 2011, pp. 1330–1335.


\bibitem{Sikora2012}
Sikora, M., Skowron, A., Wróbel, Ł. (2012). 
\newblock Rule Quality Measure-Based Induction of Unordered Sets of Regression Rules. 

\bibitem{Mann1947}
H. B. Mann and D. R. Whitney.
\newblock On a test of whether one of two random variables is stochastically larger than the other.
\newblock \emph{The Annals of Mathematical Statistics}, 18(1):50--60, 1947.

\bibitem{Suzuki1996}E. Suzuki and M. Shimura. Exceptional knowledge discovery in
databases based on information theory. In \emph{KDD'96}, pp. 275--278, 1996.


\bibitem{SuzukiUnexpected1996}E. Suzuki. Discovering unexpected exceptions:
A stochastic approach. In \emph{RSFD}, 1996.


\bibitem{Hussain2000}F. Hussain, H. Liu, E. Suzuki, and H. Lu. Exception Rule Mining
with a Relative Interestingness Measure. In \emph{PAKDD}, pp. 86--97, 2000.
DOI: 10.1007/3-540-45571-X\_11.


\bibitem{Duval2007}B. Duval, A. Salleb, and C. Vrain. On the Discovery of Exception
Rules: A Survey. In \emph{Quality Measures in Data Mining}, Studies in
Computational Intelligence, vol. 43, 2007.
DOI: 10.1007/978-3-540-44918-8\_4.


\bibitem{CalvoFlores2011}M. Calvo-Flores, M. Ruiz, and D. Sánchez. New Approaches
for Discovering Exception and Anomalous Rules. \emph{International Journal of
Uncertainty, Fuzziness and Knowledge-Based Systems}, 19:361--399, 2011.
DOI: 10.1142/S0218488511007039.


\bibitem{Vashishtha2013}J. Vashishtha, D. Kumar, and S. Ratnoo. An evolutionary
approach to discover intra-and inter-class exceptions in databases.
\emph{International Journal of Intelligent Systems Technologies and Applications},
12:283--300, 2013. DOI: 10.1504/IJISTA.2013.056535.


\bibitem{Ratnoo2018}S. Ratnoo, A. Pathak, J. Ahuja, and J. Vashishtha.
Exception Discovery Using Ant Colony Optimization.
\emph{International Journal of Computational Systems Engineering}, 4, 2018.
DOI: 10.1504/IJCSYSE.2018.10011809.


\bibitem{Wrobel2016}Ł. Wróbel, M. Sikora, and M. Michalak. Rule Quality Measures
Settings in Classification, Regression and Survival Rule Induction --- an
Empirical Approach. \emph{Fundamenta Informaticae}, 149(4):419--449, 2016.
DOI: 10.3233/FI-2016-1455.

\end{thebibliography}

\end{document}
